\documentclass[11pt]{article}

% Mandatory Packages
\usepackage{amsmath}
\usepackage{amssymb}
\usepackage{amsthm}
\usepackage{geometry}
\usepackage{tikz}
\usetikzlibrary{arrows.meta, positioning, shapes, calc}
\usepackage{mathrsfs}
\usepackage{mathtools}
\usepackage{hyperref}

% Page Geometry
\geometry{letterpaper, margin=1in}

% Theorem Environments
\newtheorem{theorem}{Theorem}[section]
\newtheorem{lemma}[theorem]{Lemma}
\newtheorem{proposition}[theorem]{Proposition}
\newtheorem{corollary}[theorem]{Corollary}
\newtheorem{definition}{Definition}[section]
\newtheorem{problem}{Problem}
\newtheorem{exercise}{Exercise}

% Custom Commands
\newcommand{\R}{\mathbb{R}}
\newcommand{\Z}{\mathbb{Z}}
\newcommand{\N}{\mathbb{N}}
\newcommand{\Q}{\mathbb{Q}}
\newcommand{\E}{\mathbb{E}}
\newcommand{\Prob}{\mathbb{P}}
\newcommand{\abs}[1]{\left|#1\right|}

\title{\textbf{Homework 4.3: Proofs Involving Real Numbers}}
\author{Tashfeen Omran}
\date{September 2026}

\begin{document}

\maketitle
\tableofcontents
\newpage

% ==========================================
% PART 1: Introduction, Context, and Prerequisites
% ==========================================
\section{Part 1: Comprehensive Introduction, Context, and Prerequisites}

\subsection{Core Concepts \& Axioms}
The study of real numbers ($\R$) in abstract mathematics transitions us from the discrete, countable domain of integers ($\Z$) to a continuous, uncountable continuum. In this section, we focus on \textbf{Proofs Involving Real Numbers}, specifically dealing with inequalities, absolute values, and foundational algebraic bounds. 

The real numbers form a \textit{complete ordered field}. This means they satisfy field axioms (closure under addition and multiplication, existence of inverses, distributivity), order axioms (trichotomy, transitivity), and the Completeness Axiom (every non-empty subset of $\R$ bounded above has a least upper bound). For our current scope, the most critical axioms and definitions are:
\begin{itemize}
    \item \textbf{Positivity Axioms:} The product and sum of two positive real numbers are positive. If $a > 0$, then $-a < 0$.
    \item \textbf{Zero Product Property:} For any $x, y \in \R$, if $xy = 0$, then $x = 0$ or $y = 0$.
    \item \textbf{Absolute Value:} The absolute value of $x \in \R$ is defined piecewise: 
    \[
    \abs{x} = 
    \begin{cases} 
      x & \text{if } x \ge 0 \\
      -x & \text{if } x < 0 
    \end{cases}
    \]
\end{itemize}

\subsection{Intuition and Logical Foundations}
When proving statements about integers, we often rely on parity (even/odd) or divisibility. Real numbers lack these discrete properties. Instead, proofs involving real numbers frequently rely on \textit{positivity} and \textit{distance}. 
The absolute value $\abs{x - y}$ intuitively represents the geometric distance between $x$ and $y$ on the real number line. Inequalities are established by showing that a difference is positive, often by manipulating an expression into a sum of squares, since $x^2 \ge 0$ for all $x \in \R$.

\subsection{Historical Context and Motivation}
In the 19th century, mathematicians like Augustin-Louis Cauchy and Karl Weierstrass realized that calculus was built on shaky, intuitive geometric foundations. To rigorously prove theorems about limits, continuity, and derivatives, they needed a formal, arithmetic definition of distance and bounds. This led to the $\epsilon-\delta$ definition of limits, which relies entirely on the manipulation of absolute values and inequalities—the exact skills developed in this chapter.

\subsection{Key Theorems and Formulas}
\begin{theorem}[The Triangle Inequality]
For every two real numbers $x$ and $y$, 
\[ \abs{x + y} \le \abs{x} + \abs{y} \]
\end{theorem}

\begin{theorem}[Arithmetic-Geometric Mean (AM-GM) Inequality]
For every two positive real numbers $a$ and $b$,
\[ \sqrt{ab} \le \frac{a + b}{2} \]
\end{theorem}

\subsection{Prerequisites}
To master this section, one must be proficient in:
\begin{enumerate}
    \item \textbf{Logical Equivalences:} Specifically, De Morgan's Laws and the contrapositive ($P \implies Q \equiv \sim Q \implies \sim P$).
    \item \textbf{Algebraic Manipulation:} Factoring polynomials, completing the square, and rationalizing denominators.
    \item \textbf{Proof by Cases:} Exhausting all possibilities (e.g., $x > 0$, $x = 0$, $x < 0$).
\end{enumerate}

\newpage
% ==========================================
% PART 2: Detailed Homework Solutions
% ==========================================
\section{Part 2: Detailed Homework Solutions}
\textit{Note: Per the strategic selection protocol, the following problems represent a comprehensive coverage of every distinct problem archetype (Equality, Contrapositive Inequalities, Absolute Values, and Means) found in Section 4.3, ensuring full conceptual mastery without redundant repetition.}

\subsection{Type 1: Algebraic Manipulation and Direct Proofs}

\begin{problem}[Exercise 4.25]
Let $x, y \in \R$. Prove that if $x^2 - 4x = y^2 - 4y$ and $x \neq y$, then $x + y = 4$.
\end{problem}
\begin{proof}
\textbf{Method: Direct Proof.} \\
Assume $x, y \in \R$ such that $x^2 - 4x = y^2 - 4y$ and $x \neq y$. We wish to show that $x + y = 4$.
We begin by rearranging the given equation to group the quadratic and linear terms:
\[ x^2 - y^2 - 4x + 4y = 0 \]
Next, we factor the difference of squares and factor out a $-4$ from the linear terms:
\[ (x - y)(x + y) - 4(x - y) = 0 \]
We can now factor out the common binomial term $(x - y)$:
\[ (x - y)(x + y - 4) = 0 \]
By the Zero Product Property of real numbers, if the product of two real numbers is zero, at least one of the factors must be zero. Therefore, either:
\[ x - y = 0 \quad \text{or} \quad x + y - 4 = 0 \]
We are given the hypothesis that $x \neq y$, which implies $x - y \neq 0$. Since the first factor cannot be zero, the second factor must be zero:
\[ x + y - 4 = 0 \]
Adding $4$ to both sides yields:
\[ x + y = 4 \]
This completes the proof.
\end{proof}

\begin{problem}[Exercise 4.28]
Prove that if $r$ is a real number such that $0 < r < 1$, then $\frac{1}{r(1-r)} \ge 4$.
\end{problem}
\begin{proof}
\textbf{Method: Direct Proof (via algebraic equivalence).} \\
Assume $r \in \R$ and $0 < r < 1$. We wish to show that $\frac{1}{r(1-r)} \ge 4$.
Since $0 < r < 1$, we know that $r > 0$ and $1 - r > 0$. Therefore, their product $r(1-r)$ is strictly positive. 
Because $r(1-r) > 0$, we can multiply both sides of the target inequality by $r(1-r)$ without changing the direction of the inequality sign. Thus, proving $\frac{1}{r(1-r)} \ge 4$ is logically equivalent to proving:
\[ 1 \ge 4r(1-r) \]
Expanding the right side, we get:
\[ 1 \ge 4r - 4r^2 \]
Rearranging all terms to the left side yields:
\[ 4r^2 - 4r + 1 \ge 0 \]
We recognize the left side as a perfect square trinomial:
\[ (2r - 1)^2 \ge 0 \]
By the properties of real numbers, the square of any real number is non-negative. Since $2r - 1 \in \R$, the inequality $(2r - 1)^2 \ge 0$ is universally true. Because our algebraic steps are reversible (since $r(1-r) > 0$), we can trace the logic backward to conclude that $\frac{1}{r(1-r)} \ge 4$.
\end{proof}

\begin{problem}[Exercise 4.36]
Prove for every positive real number $x$ that $x + \frac{1}{x} \ge 2$.
\end{problem}
\begin{proof}
\textbf{Method: Direct Proof.} \\
Assume $x \in \R$ and $x > 0$. We know that the square of any real number is non-negative. Consider the real number $(x - 1)$. We have:
\[ (x - 1)^2 \ge 0 \]
Expanding the binomial yields:
\[ x^2 - 2x + 1 \ge 0 \]
Adding $2x$ to both sides gives:
\[ x^2 + 1 \ge 2x \]
Since we assumed $x > 0$, we can divide both sides of the inequality by $x$ without changing the direction of the inequality:
\[ \frac{x^2 + 1}{x} \ge \frac{2x}{x} \]
Simplifying the fractions, we obtain:
\[ x + \frac{1}{x} \ge 2 \]
This completes the proof.
\end{proof}

\subsection{Type 2: Proof by Contrapositive}

\begin{problem}[Exercise 4.26]
Let $a, b$ and $m$ be integers. Prove that if $2a + 3b \ge 12m + 1$, then $a \ge 3m + 1$ or $b \ge 2m + 1$.
\end{problem}
\begin{proof}
\textbf{Method: Proof by Contrapositive.} \\
The statement is of the form $P \implies (Q \lor R)$. A direct proof is difficult because the conclusion is a disjunction. We will prove the contrapositive: $\sim(Q \lor R) \implies \sim P$.
By De Morgan's Laws, the negation of the conclusion is $(a < 3m + 1) \land (b < 2m + 1)$. 
Assume $a < 3m + 1$ and $b < 2m + 1$. 
Since $a, b,$ and $m$ are integers, $3m+1$ and $2m+1$ are integers. The largest integer strictly less than $3m+1$ is $3m$, and the largest integer strictly less than $2m+1$ is $2m$. Therefore, we have:
\[ a \le 3m \quad \text{and} \quad b \le 2m \]
We want to show that $\sim P$, which is $2a + 3b < 12m + 1$.
Using our inequalities, we can bound $2a$ and $3b$:
\[ 2a \le 2(3m) = 6m \]
\[ 3b \le 3(2m) = 6m \]
Adding these two inequalities together yields:
\[ 2a + 3b \le 6m + 6m = 12m \]
Since $12m < 12m + 1$, by the transitive property of inequalities, we have:
\[ 2a + 3b < 12m + 1 \]
This is exactly the negation of our original premise. Thus, the contrapositive is proven, which implies the original statement is true.
\end{proof}

\begin{problem}[Exercise 4.27]
Let $x \in \R$. Prove that if $3x^4 + 1 \le x^7 + x^3$, then $x > 0$.
\end{problem}
\begin{proof}
\textbf{Method: Proof by Contrapositive.} \\
We will prove the contrapositive: If $x \le 0$, then $3x^4 + 1 > x^7 + x^3$.
Assume $x \le 0$. We analyze the terms on the right side of the target inequality.
Since $x \le 0$, any odd power of $x$ is also less than or equal to zero. Therefore:
\[ x^7 \le 0 \quad \text{and} \quad x^3 \le 0 \]
Adding these two inequalities gives:
\[ x^7 + x^3 \le 0 \]
Now we analyze the left side of the target inequality. Since any even power of a real number is non-negative, $x^4 \ge 0$. Multiplying by 3 and adding 1 yields:
\[ 3x^4 + 1 \ge 1 > 0 \]
By transitivity, since $3x^4 + 1 > 0$ and $0 \ge x^7 + x^3$, we conclude:
\[ 3x^4 + 1 > x^7 + x^3 \]
This proves the contrapositive, establishing the truth of the original statement.
\end{proof}

\subsection{Type 3: Absolute Values and Triangle Inequality}

\begin{problem}[Exercise 4.31]
Prove for every two real numbers $x$ and $y$ that $\abs{x + y} \ge \abs{x} - \abs{y}$.
\end{problem}
\begin{proof}
\textbf{Method: Direct Proof using the Triangle Inequality.} \\
Assume $x, y \in \R$. We know by the standard Triangle Inequality that for any real numbers $a$ and $b$, $\abs{a + b} \le \abs{a} + \abs{b}$.
We can express $x$ as a sum of two terms by "adding zero" in the form of $y - y$:
\[ x = (x + y) - y \]
Taking the absolute value of both sides:
\[ \abs{x} = \abs{(x + y) + (-y)} \]
Applying the Triangle Inequality where $a = x + y$ and $b = -y$:
\[ \abs{x} \le \abs{x + y} + \abs{-y} \]
Since the absolute value of a number and its negation are equal ($\abs{-y} = \abs{y}$), we have:
\[ \abs{x} \le \abs{x + y} + \abs{y} \]
Subtracting $\abs{y}$ from both sides yields:
\[ \abs{x} - \abs{y} \le \abs{x + y} \]
Rearranging the inequality gives the desired result:
\[ \abs{x + y} \ge \abs{x} - \abs{y} \]
\end{proof}

\begin{problem}[Exercise 4.34]
Prove for every three real numbers $x, y$ and $z$ that $\abs{x - z} \le \abs{x - y} + \abs{y - z}$.
\end{problem}
\begin{proof}
\textbf{Method: Direct Proof using the Triangle Inequality.} \\
Assume $x, y, z \in \R$. We wish to bound the distance between $x$ and $z$. We can introduce $y$ into the expression by adding and subtracting $y$:
\[ x - z = x - y + y - z \]
Taking the absolute value of both sides:
\[ \abs{x - z} = \abs{(x - y) + (y - z)} \]
Let $a = x - y$ and $b = y - z$. By the Triangle Inequality ($\abs{a + b} \le \abs{a} + \abs{b}$), we obtain:
\[ \abs{(x - y) + (y - z)} \le \abs{x - y} + \abs{y - z} \]
Therefore, by substitution:
\[ \abs{x - z} \le \abs{x - y} + \abs{y - z} \]
This completes the proof.
\end{proof}

\subsection{Type 4: Means (AM-GM)}

\begin{problem}[Exercise 4.32]
(a) Recall that $\sqrt{r} > 0$ for every positive real number $r$. Prove that if $a$ and $b$ are positive real numbers, then $0 < \sqrt{ab} \le \frac{a+b}{2}$. \\
(b) Under what conditions does $\sqrt{ab} = (a+b)/2$ for positive real numbers $a$ and $b$? Justify your answer.
\end{problem}
\begin{proof}
\textbf{Part (a): Direct Proof.} \\
Assume $a$ and $b$ are positive real numbers. Since $a > 0$ and $b > 0$, their product $ab > 0$, which implies $\sqrt{ab} > 0$. This establishes the left side of the inequality ($0 < \sqrt{ab}$).
To prove the right side, consider the real numbers $\sqrt{a}$ and $\sqrt{b}$. The square of any real number is non-negative, so:
\[ (\sqrt{a} - \sqrt{b})^2 \ge 0 \]
Expanding the square yields:
\[ (\sqrt{a})^2 - 2\sqrt{a}\sqrt{b} + (\sqrt{b})^2 \ge 0 \]
Since $a$ and $b$ are positive, $(\sqrt{a})^2 = a$ and $(\sqrt{b})^2 = b$. Also, $\sqrt{a}\sqrt{b} = \sqrt{ab}$. Substituting these in:
\[ a - 2\sqrt{ab} + b \ge 0 \]
Adding $2\sqrt{ab}$ to both sides:
\[ a + b \ge 2\sqrt{ab} \]
Dividing by 2 gives:
\[ \frac{a + b}{2} \ge \sqrt{ab} \]
Combining both parts, we have $0 < \sqrt{ab} \le \frac{a+b}{2}$.

\textbf{Part (b): Condition for Equality.} \\
We want to find when $\sqrt{ab} = \frac{a+b}{2}$. Tracing our proof backward, this equality holds if and only if the initial inequality was an equality:
\[ (\sqrt{a} - \sqrt{b})^2 = 0 \]
The only real number whose square is zero is zero itself. Therefore:
\[ \sqrt{a} - \sqrt{b} = 0 \implies \sqrt{a} = \sqrt{b} \]
Squaring both sides yields $a = b$. Thus, equality holds if and only if $a = b$.
\end{proof}

\newpage
% ==========================================
% PART 3: In-Depth Analysis of Problems and Techniques
% ==========================================
\section{Part 3: In-Depth Analysis of Problems and Techniques}

\subsection{A) Problem Types and Logical Strategies}
\begin{itemize}
    \item \textbf{Algebraic Equivalence \& Non-Negativity (Exercises 4.25, 4.28, 4.36):} The overarching strategy here is to manipulate the target expression into a form involving a perfect square (e.g., $(x-1)^2 \ge 0$). Because the square of any real number is non-negative, this provides an unassailable axiomatic starting point.
    \item \textbf{Contrapositive for Inequalities (Exercises 4.26, 4.27):} When a conclusion contains an "OR" statement (disjunction) or asserts a strict positivity ($x > 0$), direct proof is often messy. The contrapositive transforms "OR" into "AND" (via De Morgan's Laws) and transforms strict inequalities into inclusive ones ($\le$), which are much easier to manipulate algebraically.
    \item \textbf{Triangle Inequality Manipulations (Exercises 4.31, 4.34):} The strategy is to use the "Adding Zero" trick. By adding and subtracting the same variable inside an absolute value (e.g., $x = x + y - y$), we can artificially create the terms needed to apply the standard Triangle Inequality $\abs{a+b} \le \abs{a} + \abs{b}$.
\end{itemize}

\subsection{B) Key Analytical and Logical Manipulations}
\begin{enumerate}
    \item \textbf{The "Adding Zero" Trick:} Used in 4.31 and 4.34. It allows us to bridge the gap between two variables by introducing an intermediate variable without changing the equation's value. This is the foundational trick for proving metric space properties.
    \item \textbf{Working Backwards (Scratchwork):} In 4.28 and 4.36, the proof starts with a seemingly random true statement like $(2r-1)^2 \ge 0$. How did we know to start there? By taking the target inequality ($\frac{1}{r(1-r)} \ge 4$) and algebraically manipulating it on scratch paper until we reached a known truth. The formal proof is then written in reverse.
    \item \textbf{Sign Analysis for Multiplication:} In 4.28, multiplying by $r(1-r)$ was only valid because we first explicitly stated that $r(1-r) > 0$. Multiplying an inequality by a variable without proving it is positive is a fatal logical flaw.
\end{enumerate}

\newpage
% ==========================================
% PART 4: "Cheatsheet", Visuals, and Tips for Success
% ==========================================
\section{Part 4: "Cheatsheet", Visuals, and Tips for Success}

\subsection{Theorems \& Equivalences}
\begin{itemize}
    \item \textbf{Contrapositive:} $P \implies Q \equiv \sim Q \implies \sim P$
    \item \textbf{De Morgan's Laws:} $\sim(P \lor Q) \equiv \sim P \land \sim Q$ \quad and \quad $\sim(P \land Q) \equiv \sim P \lor \sim Q$
    \item \textbf{Triangle Inequality:} $\abs{x + y} \le \abs{x} + \abs{y}$
    \item \textbf{Reverse Triangle Inequality:} $\abs{x - y} \ge \abs{\abs{x} - \abs{y}}$
    \item \textbf{AM-GM Inequality:} $\frac{x_1 + x_2 + \dots + x_n}{n} \ge \sqrt[n]{x_1 x_2 \dots x_n}$
\end{itemize}

\subsection{Set \& Logic Illustrations: Geometric Proof of AM-GM}
The Arithmetic-Geometric Mean inequality has a beautiful geometric interpretation. Consider a semicircle with diameter $a+b$. The radius is the Arithmetic Mean. A perpendicular chord segment from the diameter to the circle represents the Geometric Mean. Since a chord segment can never exceed the radius, $\sqrt{ab} \le \frac{a+b}{2}$.

\begin{center}
\begin{tikzpicture}[scale=1.5, thick]
    % Define variables a and b for the drawing
    \def\a{4}
    \def\b{2}
    \def\rad{(\a+\b)/2} % Radius = 3
    
    % Draw the semicircle
    \draw[blue, thick] (\a+\b, 0) arc (0:180:\rad);
    
    % Draw the diameter line
    \draw[black, thick] (0,0) -- (\a+\b,0);
    
    % Draw the center point
    \coordinate (C) at (\rad, 0);
    \filldraw[black] (C) circle (1pt) node[below] {Center};
    
    % Draw the junction point
    \coordinate (D) at (\a, 0);
    \filldraw[black] (D) circle (1pt);
    
    % Calculate height of geometric mean
    \pgfmathsetmacro{\gm}{sqrt(\a*\b)}
    \coordinate (E) at (\a, \gm);
    
    % Draw Geometric Mean line
    \draw[red, thick] (D) -- (E) node[midway, right] {$\sqrt{ab}$ (GM)};
    
    % Draw Radius line (Arithmetic Mean)
    \draw[green!60!black, thick] (C) -- (E) node[midway, left=2pt] {$\frac{a+b}{2}$ (AM)};
    
    % Labels for a and b
    \draw[<->] (0, -0.3) -- (\a, -0.3) node[midway, below] {$a$};
    \draw[<->] (\a, -0.3) -- (\a+\b, -0.3) node[midway, below] {$b$};
    
    % Right angle square
    \draw (\a-0.15, 0) -- (\a-0.15, 0.15) -- (\a, 0.15);
\end{tikzpicture}
\end{center}

\subsection{Heuristics \& "Cheats"}
\begin{itemize}
    \item \textbf{The "OR" Trigger:} If the conclusion of an implication contains an "OR", immediately try a proof by contrapositive.
    \item \textbf{The "Square" Trigger:} If you need to prove an expression is $\ge 0$, try to factor it into a perfect square or a sum of perfect squares.
    \item \textbf{The "Absolute Value" Trigger:} If a problem asks to bound an absolute value $\abs{A - B}$, try introducing a middle term $C$ and use the Triangle Inequality: $\abs{A - B} \le \abs{A - C} + \abs{C - B}$.
\end{itemize}

\subsection{Pitfalls \& Fallacies}
\begin{itemize}
    \item \textbf{Dividing by an Unknown Sign:} Dividing an inequality by $x$ without first proving $x > 0$. If $x < 0$, the inequality sign must flip.
    \item \textbf{Assuming the Consequent (Circular Reasoning):} Starting a proof by writing down the inequality you are trying to prove, manipulating it to $0 = 0$, and declaring victory. You must start from a known truth and work \textit{forward} to the target.
\end{itemize}

\newpage
% ==========================================
% PART 5: Conceptual Synthesis and The "Big Picture"
% ==========================================
\section{Part 5: Conceptual Synthesis and The "Big Picture"}

\subsection{A) Thematic Connections}
The core theme of this section is \textbf{formalizing distance and bounds in a continuous space}. In earlier chapters, proofs relied on discrete properties (an integer is either even or odd). In $\R$, we rely on the topology of the real line. The Triangle Inequality is not just an algebraic trick; it is the defining axiom of a \textit{Metric Space}. It formalizes the geometric intuition that "the shortest path between two points is a straight line."

\subsection{B) Forward and Backward Links}
\begin{itemize}
    \item \textbf{Backward Links:} This section builds directly on the logical frameworks (Direct Proof, Contrapositive) established in Chapters 1-3, applying them to the axioms of real numbers.
    \item \textbf{Forward Links:} The inequalities mastered here are the absolute bedrock of \textbf{Real Analysis}. The $\epsilon-\delta$ proofs for limits and continuity consist entirely of manipulating absolute values using the Triangle Inequality. Furthermore, the AM-GM inequality is a gateway to \textbf{Convex Optimization} and Jensen's Inequality, which are foundational in Machine Learning and Probability Theory.
\end{itemize}

% ==========================================
% PART 6: Abstract Real-World Modeling
% ==========================================
\section{Part 6: Abstract Real-World Modeling (AI, ML, \& Quant Finance)}

\subsection{A) Concrete Scenarios}
\begin{enumerate}
    \item \textbf{Machine Learning (L1 Regularization / Lasso):} In training neural networks, we often want sparse weights (many weights equal to exactly zero) to prevent overfitting. We add an L1 penalty term to the loss function: $\lambda \sum \abs{w_i}$. Proving that this penalty space is convex and bounds the parameter space relies entirely on the Triangle Inequality ($\abs{w_1 + w_2} \le \abs{w_1} + \abs{w_2}$).
    \item \textbf{Quantitative Finance (Coherent Risk Measures):} When a bank calculates the risk of a portfolio, the risk measure $\rho(X)$ must be "coherent". One of the axioms of coherence is \textit{Subadditivity}: $\rho(X + Y) \le \rho(X) + \rho(Y)$. This ensures that diversification reduces risk. This is a direct abstraction of the Triangle Inequality.
\end{enumerate}

\subsection{B) Model Problem Setup: Coherent Risk in Portfolio Optimization}
Let $(\Omega, \mathcal{F}, \Prob)$ be a probability space representing all possible market states. Let $X, Y \in L^2(\Omega, \mathcal{F}, \Prob)$ be random variables representing the future returns of two financial portfolios. 

We define a risk measure $\rho: L^2 \to \R$. For $\rho$ to be a \textit{Coherent Risk Measure} (Artzner et al., 1999), it must satisfy Subadditivity:
\[ \rho(X + Y) \le \rho(X) + \rho(Y) \quad \forall X, Y \in L^2 \]
Suppose we define our risk measure using the $L^2$ norm (Standard Deviation of losses), such that $\rho(X) = \sqrt{\E[X^2]}$. 
To prove this risk measure is coherent, we must prove the stochastic equivalent of the Triangle Inequality (Minkowski's Inequality for $p=2$):
\[ \sqrt{\E[(X+Y)^2]} \le \sqrt{\E[X^2]} + \sqrt{\E[Y^2]} \]
\textbf{Optimization Problem:} A quantitative trader wishes to find portfolio weights $w_1, w_2$ such that $w_1 + w_2 = 1$ to minimize total risk $\rho(w_1 X + w_2 Y)$ subject to a minimum expected return constraint $\E[w_1 X + w_2 Y] \ge \mu_{target}$. The subadditivity property guarantees that the risk surface is convex, ensuring that gradient descent algorithms will converge to a unique global minimum rather than getting stuck in local optima.

\newpage
% ==========================================
% PART 7: Common Variations and Untested Concepts
% ==========================================
\section{Part 7: Common Variations and Untested Concepts}

A standard advanced topic missing from the provided homework is the \textbf{Cauchy-Schwarz Inequality} for real numbers. It is the natural extension of the AM-GM inequality and is fundamental in linear algebra, probability (bounding covariance), and machine learning (cosine similarity).

\begin{theorem}[Cauchy-Schwarz Inequality in $\R^2$]
For any real numbers $a, b, x, y \in \R$,
\[ (ax + by)^2 \le (a^2 + b^2)(x^2 + y^2) \]
\end{theorem}
\begin{proof}
Consider the function $f(t) = (at - x)^2 + (bt - y)^2$ for a real variable $t$.
Since it is a sum of squares, $f(t) \ge 0$ for all $t \in \R$.
Expanding the expression, we get:
\[ f(t) = a^2t^2 - 2axt + x^2 + b^2t^2 - 2byt + y^2 \]
Grouping the terms by powers of $t$:
\[ f(t) = (a^2 + b^2)t^2 - 2(ax + by)t + (x^2 + y^2) \]
This is a quadratic equation in $t$ of the form $At^2 + Bt + C \ge 0$, where $A = a^2+b^2$, $B = -2(ax+by)$, and $C = x^2+y^2$.
For a quadratic to be non-negative for all real $t$, it cannot have two distinct real roots. Therefore, its discriminant must be less than or equal to zero ($\Delta \le 0$):
\[ B^2 - 4AC \le 0 \]
Substituting our values for $A, B,$ and $C$:
\[ (-2(ax + by))^2 - 4(a^2 + b^2)(x^2 + y^2) \le 0 \]
\[ 4(ax + by)^2 - 4(a^2 + b^2)(x^2 + y^2) \le 0 \]
Dividing by 4 and rearranging yields the Cauchy-Schwarz Inequality:
\[ (ax + by)^2 \le (a^2 + b^2)(x^2 + y^2) \]
\end{proof}

\newpage
% ==========================================
% PART 8: Advanced Diagnostic Testing: "Find the Flaw"
% ==========================================
\section{Part 8: Advanced Diagnostic Testing: "Find the Flaw"}
\textit{Act as the grader. Each of the following proofs contains a single, critical logical or mathematical flaw. Identify the error, explain why it is invalid, and provide the correct logical step.}

\subsection{Flawed Proofs}

\begin{problem}
Prove that for all real numbers $x, y$, if $x^2 > y^2$, then $x > y$.
\end{problem}
\begin{proof}[Flawed Proof 1]
Assume $x^2 > y^2$. Taking the square root of both sides yields $\sqrt{x^2} > \sqrt{y^2}$. Since the square root and the square cancel each other out, we are left with $x > y$. This completes the proof.
\end{proof}

\begin{problem}
Prove that for all positive real numbers $x$, $x + \frac{4}{x} \ge 4$.
\end{problem}
\begin{proof}[Flawed Proof 2]
We want to show that $x + \frac{4}{x} \ge 4$. 
Multiply both sides by $x$ to get $x^2 + 4 \ge 4x$.
Subtract $4x$ from both sides to get $x^2 - 4x + 4 \ge 0$.
Factor the left side to get $(x - 2)^2 \ge 0$.
Since the square of any real number is non-negative, $(x - 2)^2 \ge 0$ is true. Therefore, the original statement is true.
\end{proof}

\begin{problem}
Prove that for any real numbers $x, y$, $\abs{x - y} \le \abs{x} - \abs{y}$.
\end{problem}
\begin{proof}[Flawed Proof 3]
By the Triangle Inequality, we know that $\abs{a + b} \le \abs{a} + \abs{b}$. 
Let $a = x$ and $b = -y$. Substituting these into the Triangle Inequality gives:
$\abs{x + (-y)} \le \abs{x} + \abs{-y}$.
Since $\abs{-y} = -\abs{y}$, we substitute this in to get:
$\abs{x - y} \le \abs{x} - \abs{y}$.
\end{proof}

\begin{problem}
Prove that if $x \in \R$ and $x^3 < 0$, then $x < 0$.
\end{problem}
\begin{proof}[Flawed Proof 4]
We will use a proof by contrapositive. The contrapositive of the statement is: "If $x > 0$, then $x^3 > 0$."
Assume $x > 0$. Multiplying the inequality $x > 0$ by $x^2$ (which is positive) gives $x^3 > 0$. 
Since the contrapositive is true, the original statement is true.
\end{proof}

\begin{problem}
Prove that for any real numbers $a, b$, if $ab = 0$, then $a = 0$ and $b = 0$.
\end{problem}
\begin{proof}[Flawed Proof 5]
Assume $ab = 0$. We will prove this by cases.
Case 1: Assume $a = 0$. Then $0 \cdot b = 0$, which is true.
Case 2: Assume $b = 0$. Then $a \cdot 0 = 0$, which is true.
Since both cases result in $0 = 0$, it must be that $a = 0$ and $b = 0$.
\end{proof}

\newpage
\subsection{Solution Key: Flaw Resolutions}

\textbf{Resolution to Flawed Proof 1:}
\begin{itemize}
    \item \textbf{The Flaw:} The assertion that $\sqrt{x^2} = x$ is false for negative numbers. The correct identity is $\sqrt{x^2} = \abs{x}$.
    \item \textbf{Why it's an error:} If $x = -3$ and $y = 2$, then $(-3)^2 > 2^2$ ($9 > 4$), but $-3$ is not greater than $2$.
    \item \textbf{Correction:} The statement itself is false. It should be corrected to state: "If $x^2 > y^2$, then $\abs{x} > \abs{y}$."
\end{itemize}

\textbf{Resolution to Flawed Proof 2:}
\begin{itemize}
    \item \textbf{The Flaw:} Circular reasoning / Assuming the consequent. The proof starts with the statement it is trying to prove and works backward to a known truth.
    \item \textbf{Why it's an error:} Logical implications do not automatically work in reverse. Just because $P \implies Q$ and $Q$ is true, does not mean $P$ is true (Affirming the Consequent).
    \item \textbf{Correction:} The proof must be written in reverse. Start with the known truth: "Since $x \in \R$, $(x-2)^2 \ge 0$." Then expand to $x^2 - 4x + 4 \ge 0$, rearrange to $x^2 + 4 \ge 4x$, and \textit{then} state "Since $x > 0$, we can divide by $x$ to get $x + \frac{4}{x} \ge 4$."
\end{itemize}

\textbf{Resolution to Flawed Proof 3:}
\begin{itemize}
    \item \textbf{The Flaw:} The algebraic substitution $\abs{-y} = -\abs{y}$ is completely false.
    \item \textbf{Why it's an error:} Absolute value is always non-negative. $\abs{-y} = \abs{y}$, not $-\abs{y}$. For example, if $y = 3$, $\abs{-3} = 3$, but $-\abs{3} = -3$.
    \item \textbf{Correction:} The correct application of the Triangle Inequality yields $\abs{x - y} \le \abs{x} + \abs{y}$. The statement $\abs{x - y} \le \abs{x} - \abs{y}$ is mathematically false (e.g., $x=0, y=1 \implies 1 \le -1$, false).
\end{itemize}

\textbf{Resolution to Flawed Proof 4:}
\begin{itemize}
    \item \textbf{The Flaw:} The contrapositive was formed incorrectly.
    \item \textbf{Why it's an error:} The negation of $x < 0$ is $x \ge 0$, not $x > 0$. Similarly, the negation of $x^3 < 0$ is $x^3 \ge 0$.
    \item \textbf{Correction:} The correct contrapositive is: "If $x \ge 0$, then $x^3 \ge 0$." The proof should then proceed: Assume $x \ge 0$. If $x = 0$, $x^3 = 0 \ge 0$. If $x > 0$, $x^3 > 0 \ge 0$. Thus $x^3 \ge 0$.
\end{itemize}

\textbf{Resolution to Flawed Proof 5:}
\begin{itemize}
    \item \textbf{The Flaw:} The proof confuses the logical "OR" with the logical "AND", and it attempts to prove an implication by assuming the conclusion.
    \item \textbf{Why it's an error:} The Zero Product Property states that if $ab = 0$, then $a = 0$ \textbf{OR} $b = 0$. The statement "$a=0$ \textbf{AND} $b=0$" is false (e.g., $a=5, b=0$ gives $ab=0$). Furthermore, plugging $a=0$ into the equation to get $0=0$ does not prove that $a$ \textit{must} be $0$.
    \item \textbf{Correction:} The original statement is false. The correct theorem is $ab = 0 \implies a = 0 \lor b = 0$. The proof should assume $ab = 0$ and $a \neq 0$, then multiply by $1/a$ to show $b = 0$.
\end{itemize}

\end{document}